\documentclass[10pt,a4paper]{amsart}
\usepackage[margin=19mm]{geometry}
\usepackage{amsmath,amssymb,mathtools}
\usepackage[scr=rsfs,cal=euler]{mathalfa}
\usepackage{xfrac}
\usepackage[unicode,hidelinks,bookmarks=false]{hyperref}
\newcommand{\ie}{i.e.\ }
\newcommand{\cf}{cf.\ }
\newcommand{\resp}{respectively\ }
\newcommand{\Subsection}{\subsection}
\providecommand{\nobreakdash}{\nobreak\hbox{-}}
\setlength{\emergencystretch}{2em}
\multlinegap0pt
\usepackage{bm}
\usepackage{bbm}
\usepackage{dsfont}
\usepackage{xspace}
\usepackage{cancel}
\usepackage{mathtools}
\usepackage{cleveref}
\usepackage{amsthm}
\makeatletter
\@ifundefined{thm}{\newtheorem{thm}{Thm}}{}
\@ifundefined{lem}{\newtheorem{lem}[thm]{Lemma}}{}
\makeatother
\newcommand{\beq}[1]{\begin{equation}\label{#1}}
\newcommand{\cG}{\mathcal{G} }
\newcommand{\cV}{\mathcal{V} }
\newcommand{\cW}{\mathcal{W} }
\newcommand{\enq}[0]{\end{equation}}
\newcommand{\nin}[0]{\noindent}
\newcommand{\sub}[0]{\subseteq}
\title[A refined graph container lemma and applications to the hard]{A refined graph container lemma and applications to the hard-core model on bipartite expanders}
\author{Matthew Jenssen, Alexandru Malekshahian, Jinyoung Park}
\date{Source: arXiv:2411.03393v3 (CC BY 4.0)}
\begin{document}
\maketitle

\noindent\emph{Source and licence.} A refined graph container lemma and applications to the hard-core model on bipartite expanders. Version arXiv:2411.03393v3 (CC BY 4.0), distributed under the Creative Commons Attribution 4.0 International licence, as recorded in the arXiv licence metadata for this exact version and shown on the abstract page https://arxiv.org/abs/2411.03393v3. Full rights notice: licenses/arxiv-2411.03393-CC-BY-4.0.txt.\par\smallskip

\noindent\textit{Editorial scope.} Counted unit: the Lem whose source label is \texttt{lem.heart} in the article (source0.tex lines 493--497, source label \texttt{lem.heart}), delivered together with the article's own complete original proof of it (source0.tex lines 499--517, one \texttt{proof} environment of 19 lines). No mathematics is written, repaired or re-derived: statement and proof are the article's own text line for line. Required context retained uncounted: lem.psi (lines 463--466). Every definition, display and label of those lines is retained unchanged. Omitted from this package and not redistributed: 1--462, 467--492, 498--498, 518--920. Nothing inside a retained excerpt is omitted or altered: every retained body is delivered exactly as the article's lines are written, so no omission receipt of any kind accompanies this package. Citations: the retained excerpts carry no citation command at all, so no bibliography entry of the article is transcribed and no reference is redistributed. Reference adaptation: none was needed and none was made; every reference inside the delivered unit resolves inside it, so no unresolved reference or citation is produced. Printed slips kept literal: no printed word, symbol, hypothesis or number of the counted statement or of its proof was altered. Layout and class: the article's own class is \texttt{amsart} with packages including amsmath, amssymb, amsbsy, array, graphpap, color, paralist, pstricks, eucal, graphicx, hyperref, pifont, mathtools, cleveref; the wrapper is the lane's pinned \texttt{amsart} editorial wrapper at 10pt on A4, which restates the theorem-like environments the retained text needs and adds the article's own macro definitions that the retained text needs (\texttt{\textbackslash beq}, \texttt{\textbackslash cG}, \texttt{\textbackslash cV}, \texttt{\textbackslash cW}, \texttt{\textbackslash enq}, \texttt{\textbackslash nin}, \texttt{\textbackslash sub}), with extra wrapper packages (bm, bbm, dsfont, xspace, cancel, mathtools, cleveref). The delivered numbering is the unit's own. Assets: no retained span contains a figure environment, a graphics-inclusion command, a PDF-page inclusion, a table or a TikZ picture; no image file is copied into this package, the package declares no asset dependency and no image file is redistributed with it. Licence: version arXiv:2411.03393v3 of this article is distributed under the Creative Commons Attribution 4.0 International licence, as recorded in the arXiv licence metadata for this exact version and shown on the abstract page https://arxiv.org/abs/2411.03393v3. A full rights notice is delivered at licenses/arxiv-2411.03393-CC-BY-4.0.txt. Characters: every retained excerpt is pure ASCII, so the delivered engine, XeLaTeX with its default Latin Modern text font, reports no missing character.
\begin{lem}[Cost for initial run]\label{lem.psi} For any $F' \in \cV(a,g,\varphi)$ and $1 \le \psi_X \leq d_X-1, \ 1\leq \psi_Y \le d_Y/\delta-1$ there is a family $\cW=\cW(F', \psi_X, \psi_Y) \sub 2^Y \times 2^X$ with
\beq{psi_bd}|\cW|\le {d_Yg \choose \le \delta w/((d_Y-\delta \varphi)\psi_X)}{\delta d_X d_Y g \choose \le w/((d_X-\psi_X)\psi_Y)}\enq
such that any $A \in \cG(a,g)$ for which $F'$ is a $\varphi$-approximation has a $(\psi_X, \psi_Y)$-approximating pair in $\cW.$
\end{lem}

\begin{lem}[Cost for additional run]\label{lem.heart} Let $\cW=\cW(F', \psi_X, \psi_Y)$ be as in \Cref{lem.psi}.
    For each $(F,S) \in \cW$, if $1 \le \psi_X'\leq \psi_X, \ 1\leq \psi_Y' \le \psi_Y$, then there is a family $\cW'=\cW'(F,S) \sub 2^Y \times 2^X$ with
    \beq{psi'_bd}|\cW'|\le {|S| \choose \le (g-|F|)/\psi_X'}{\delta d_X|S| \choose \le w/((d_X-\psi_X')\psi_Y')}\enq
    such that any $A \in \cG(a,g)$ for which $(F,S)\in \cW$ is a $(\psi_X, \psi_Y)$-approximating pair has a $(\psi_X', \psi_Y')$-approximating pair in $\cW'.$
\end{lem}

\begin{proof}
    For $(F,S)$ as in the statement of the lemma, we additionally run the algorithm in the proof of \Cref{lem.psi} with $\psi_X'$ and $\psi_Y'$.% Note that $\psi_X' = \psi_X \cdot O(1/K)=O(\gamma \psi_X)$, since $s-a=\gO(t)$ and $g-f\stackrel{\eqref{eq.SF}}{=}O(t)$ (\eqref{eq.SF} implies that $g-f \le (1-\psi_Y/d)^{-1}(g-s+(s-a)\psi_X/d)=O(t)$). 
\begin{quote}
    \nin \textbf{Input.} $A \in \cG(a,g)$ and its $(\psi_X, \psi_Y)$-approximating pair $(F,S)$

    \nin \textbf{Step 1'.}  If $\{u \in [A]:d_{N(A)\setminus F}(u)>\psi_X'\} \ne \emptyset,$ pick the smallest $u$ (in $\prec$) in this set and update $F$ by $F \leftarrow F \cup N(u).$ Repeat this until $\{u \in [A]:d_{N(A) \setminus F}(u)>\psi_X'\}=\emptyset.$ Then set $\hat F=F$ and $\hat S=\{u \in S:d_{\hat F}(u) \ge d(u)-\psi_X'\}$ and go to Step 2'.

    \nin \textbf{Step 2'.} If $\{v \in Y \setminus N(A): d_{\hat S}(v)>\psi_Y'\} \ne \emptyset$, pick the smallest $v$ (in $\prec$) in this set and update $\hat S$ by $\hat S \leftarrow \hat S \setminus N(v).$ Repeat this until $\{v \in Y \setminus N(A):d_{\hat S}(v)>\psi_Y'\}=\emptyset.$ Then set $\tilde S=\hat S$ and $\tilde F=\hat F \cup \{v \in Y:d_{\tilde S}(v)>\psi_Y'\}$ and stop.

    \nin \textbf{Output.} $(\tilde F, \tilde S)$
\end{quote}
It is easy to see that $(\tilde F, \tilde S)$ is a $(\psi_X', \psi_Y')$-approximating pair for $A \in \cG(a,g)$. We bound the number of outputs for each $(F,S)$. Write $f=|F|.$

    \nin \textit{Cost for Step 1'.} Each iteration in Step 1' removes at least $\psi_X'$ vertices from $N(A) \setminus F$, a set of size $g-f$, and so there can be at most ${(g-f)}/{\psi_X'}$ iterations. The $u$'s in Step 1' are all drawn from $[A]$ and hence from $S$. So the total number of outputs for Step 1' is at most \[{|S| \choose \le (g-f)/\psi_X'}.\]

    \nin \textit{Cost for Step 2'.} Each $u \in \hat S \setminus [A]$ contributes more than $d_X-\psi_X'$ edges to $\nabla(N(A), X \setminus [A])$, so initially $|\hat S \setminus [A]| \le w/(d_X-\psi_X').$ Each $v$ used in Step 2' reduces this by at least $\psi_Y',$ so there are at most $w/((d_X-\psi_X')\psi_Y')$ iterations. Each $v$ is drawn from $N(\hat S),$ a set which is contained in $N(S)$ and so has size at most $\delta d_X|S|.$ So the total number of outputs for Step 2' is at most
    \[{\delta d_X|S| \choose \le w/((d_X-\psi_X')\psi_Y')}. \qedhere\]

\end{proof}

\end{document}
